\section{Related Work}
\subsection{Privacy and trustworthy crowdsensing in vehicular networks}
Research on vehicular crowdsensing security has addressed participant identity, confidentiality of sensed data, and evidence trustworthiness. AARE combines anonymous authentication with privacy-preserving assessment of data reliability~\citep{Cheng2023}. BARM introduces blockchain-assisted anonymous authentication and reputation management to evaluate and update participant reputations~\citep{Zheng2025BARM}. Regarding data processing, PPVDA studies privacy-preserving and verifiable aggregation~\citep{Zhou2024}; BCS-LPP employs techniques including Geohash and order-preserving encryption to protect sensing locations while supporting location checks and task management~\citep{Shen2024}. Cooperative detection of anomalous vehicular behavior has also been investigated at both vehicle and roadside-unit levels~\citep{Sultana2024}. Other recent work examines lightweight anonymous authentication, certificateless aggregate signatures, and SNARK-based verification of computations on authenticated data~\citep{Imam2025JISA,Wang2025JISA,Reinhart2026ADSC}. Recent VANET work in the \emph{Journal of Network and Computer Applications} also studies conditional privacy-preserving message authentication with secret-key updates~\citep{JNCaFan2026}; its authentication objective is distinct from proving the spatial compliance of mutually uncertain signed reports.

At the vehicular edge, Li et al. study fog-based IoV security using post-quantum signcryption and blockchain-supported trust management~\citep{JNCaLi2025}. Their work addresses secure communications and dynamic trust, while our module concentrates on the semantics of authenticated spacetime uncertainty.

Authentication has also been studied in other adversarial network settings. Umrani and Palmieri examine threshold password-authenticated key exchange for vehicular ad hoc networks~\citep{CybersecUmrani2025}, while Wang et al. develop multi-authority aggregate signatures for privacy-preserving message authentication in smart grids~\citep{CybersecWang2026}. Their security objectives concern key establishment and message authenticity, not whether signed measurement regions physically cover a reported event.

In mobile crowdsensing, the BLIND system protects data sources and improves truth discovery through encrypted clustering and private outlier handling~\citep{JNCaAgate2024}. It addresses data-quality inference rather than a report-bound certificate that all sufficiently supported spacetime states meet a public policy.

These mechanisms address different parts of the sensing workflow: authentication governs participation, aggregation protects inputs, and misbehavior detection identifies suspicious participants. None of these functions by itself resolves whether an event complies with a policy when authenticated sources provide uncertain or conflicting observations. That decision also needs to account for uncertainty in the individual observations.

\subsection{Zero-knowledge location and geometric-policy verification}
Location privacy can be pursued through application-specific obfuscation and protected queries~\citep{Mdee2022,Atmaca2024}, or by credentials and authentication protocols that limit the disclosure of sensitive location information~\citep{Ghani2025,Zhang2025Sybil,Wang2025ZKAuth}. Zero-knowledge proofs offer another approach: a verifier can check a geometric constraint without learning the protected coordinate or an entire trajectory. ZKLP, proposed by Ernstberger and colleagues, incorporates floating-point geometry into zero-knowledge proofs for region membership and proximity predicates~\citep{ZKLP}. The vehicle-policy protocols of Bogdanov and colleagues instead use locations signed by a trusted witness to prove compliance with driving-area requirements in subsidy or taxation applications~\citep{VehiclePolicy}.

Privacy-preserving policy enforcement has also been examined outside geometric verification. Kerl et al. investigate private attribute-based access decisions using homomorphic encryption and zero-knowledge proofs~\citep{CybersecKerl2025}, while Zhang develops zero-knowledge policy checks over encrypted network traffic~\citep{CybersecZhang2026}. These studies certify attribute or traffic-policy relations rather than threshold support across signed, uncertain spacetime observations.

The location-proof systems discussed above demonstrate that geometric predicates and vehicle-trajectory policies can be certified while withholding the underlying positions. Protecting a position, however, differs from establishing that the asserted position is adequately supported by observations. With disagreeing measurement intervals or sources that may maliciously sign data, proving membership of one chosen point or an existing trajectory does not establish that \emph{all} states supported by sufficient evidence satisfy the policy. Our proof relation therefore concerns the threshold support induced jointly by multiple signed intervals, rather than repeated membership testing of a selected coordinate.

\subsection{Byzantine-robust sensor fusion}
Uncertainty and disagreement across measurements have long been studied in fault-tolerant sensor fusion. Marzullo's approach to continuous-valued sensor failures uses overlapping intervals to characterize states consistent with potentially faulty observations~\citep{Marzullo1990}. BFT-PoLoc applies Byzantine-resistant geometry to localization from Internet delays~\citep{BFTPoLoc}. PG combines private computation and Byzantine-tolerant sensor fusion while addressing output delivery~\citep{PG}. CRSF further examines private fusion when sensors and servers may collude, including participation consistency and protocol liveness~\citep{CRSF}.

These systems handle inconsistent observations and malicious participants; some also provide participation or delivery guarantees beyond those of a fixed issuer roster. Their output is a fused estimate or a secure computation result. Our setting instead calls for a report-bound proof that a verifier can check against a spacetime policy without seeing the signed intervals.

\subsection{Comparison and research positioning}
Table~\ref{tab:related-work} compares the objects verified by these methods, their trust assumptions, and their outputs. Our construction uses threshold support to retain the uncertainty in signed evidence, then proves the resulting set-level policy decision in zero knowledge.

\begin{table}[htbp]\centering\footnotesize
\caption{Functional scope and trust assumptions of representative systems.}\label{tab:related-work}
\begin{tabularx}{\linewidth}{@{}p{0.18\linewidth}p{0.21\linewidth}p{0.29\linewidth}X@{}}
\toprule Method & Object & Evidence and trust assumptions & Privacy and output\\\midrule
BARM~\citep{Zheng2025BARM} & Identity and reputation & Reputation evaluation and blockchain-assisted identity management & Anonymous authentication and private reputation evaluation\\
PPVDA~\citep{Zhou2024} & Sensing aggregates & Homomorphic message authentication, secret sharing, malicious participants & Private inputs and verifiable aggregation\\
BCS-LPP~\citep{Shen2024} & Sensing location & Geohash, order-preserving encryption, deposit mechanism & Location privacy and task/location checking\\
BFT-PoLoc~\citep{BFTPoLoc} & Delay-based location & Internet delays and Byzantine-resistant geometry & Ranging and multilateration; privacy not rated here\\
PG~\citep{PG} & Fused sensor result & Malicious sensors and aggregating servers & Private inputs/results; output-delivery guarantee\\
CRSF~\citep{CRSF} & Fusion and participation & Sensor--server collusion and consistency of submissions & Private fusion, explicit abort and liveness\\
ZKLP~\citep{ZKLP} & Region and proximity & Hidden coordinates; floating-point geometry & Zero-knowledge geometric proof\\
Vehicle policy~\citep{VehiclePolicy} & Driving-area policy & Locations signed by a trusted witness & Hidden trajectory; subsidy and tax compliance\\
\textbf{This work} & \textbf{All threshold-supported spacetime states} & \textbf{Signed uncertainty boxes, fixed roster, budgets $f$ and $r$} & \textbf{Hidden boxes/signatures; report-bound four-valued proof}\\\bottomrule
\end{tabularx}
\par\smallskip\raggedright\footnotesize BFT-PoLoc and CRSF refer to the cited arXiv versions. All privacy and correctness claims depend on their respective trust assumptions. Semantic inconsistency must not be confused with abort on missing evidence or proof failure.
\end{table}

The circuit checks issuer signatures and set-level compliance for a pre-established roster. It produces a compact, noninteractive certificate bound to the task, report, and challenge; the physical interpretation remains conditional on the fault and coverage budgets.
